\subsection{معرفی و انگیزه}
این مکانیزم بر مبنای مبانی نظری پیشرفته در ادبیات توجه و فشرده‌سازی \cite{ref_flap_2024} توسعه یافته است.

هرس ساختاریافته سنتی در مدل‌های ترنسفورمر اغلب با اتکا بر وزن‌های ایستا (مانند نُرم $\ell_1$ یا $\ell_2$) انجام می‌شود که بدون در نظر گرفتن دینامیک فعال‌سازی‌های ورودی، ستون‌ها یا کانال‌هایی را حذف می‌کند که در واقعیت برای حفظ بازنمایی داده‌ها حیاتی هستند. از سوی دیگر، روش‌های متکی بر اطلاعات مرتبه دوم، پیچیدگی محاسباتی بالایی داشته و سربار زمانی قابل توجهی در پردازش شبکه‌های بزرگ ایجاد می‌کنند.

ماژول فشرده‌سازی \lr{FLAP (Fluctuation-Based Adaptive Structured Pruning)} به عنوان یک راهکار هرس ساختاریافته تطبیقی مبتنی بر نوسانات فعال‌سازی پدید آمده است. ایده محوری این روش، بهره‌گیری از معیار واریانس فعال‌سازی‌های ورودی (\lr{Activation Fluctuation/Variance}) به عنوان عاملی کلیدی در ترکیب با انرژی وزن‌ها است. کانال‌ها یا ابعادی از شبکه که فعال‌سازی آن‌ها در داده‌های کالیبراسیون دارای نوسان بالاتری هستند، توانمندی بیشتری در تمایز و تفکیک توکن‌ها دارا بوده و نباید هرس شوند. این ماژول با پشتیبانی بومی از مدل‌های انکودر و دیکودر، امکان هرس کانال و همچنین تنکی ساختاریافته $2:4$ (\lr{N:M Sparsity}) شتاب‌یافته روی پردازنده‌های گرافیکی مدرن را فراهم می‌سازد.
